\documentclass{article}

\usepackage[final]{neurips_2024}
\usepackage[utf8]{inputenc}
\usepackage[T1]{fontenc}
\usepackage{hyperref}
\usepackage{url}
\usepackage{booktabs}
\usepackage{amsfonts}
\usepackage{amssymb}
\usepackage{nicefrac}
\usepackage{microtype}
\usepackage{xcolor}
\usepackage{graphicx}

\title{Liquidity-weighted Polymarket vs.\ Michigan Consumer Sentiment:\\ a 2022--2025 backtest}

\author{
  Aly Sultan \\
  Independent Researcher \\
  \texttt{aly.sultan@example.com}
  \And
  Mira Hosseini \\
  Department of Economics \\
  Northbridge University \\
  \texttt{m.hosseini@northbridge.edu}
  \And
  Tomas \"O\v zek \\
  Bratislava School of Public Policy \\
  \texttt{t.ozek@bspp.sk}
}

\begin{document}
\maketitle

\begin{abstract}
We construct a daily, liquidity-weighted macro-pessimism index from $1{,}107$
resolved Polymarket contracts on US recession, inflation, Federal Reserve
actions, and labor-market outcomes, and compare it against the monthly
University of Michigan Index of Consumer Sentiment (ICS) over the 39 months
during which the two series overlap (Feb 2023 -- Apr 2026). The aggregate
index is essentially uncorrelated with ICS in both levels
($r = +0.09,\ p = 0.60$) and first differences
($r = -0.03,\ p = 0.87,\ R^2 = 0.0007$). The result is not an artefact of
the equal-weighting scheme (equal-weighted: $r_{\Delta} = +0.32$, $p = 0.05$,
but with the \emph{wrong} sign for our hypothesis) and survives restriction
to high-volume markets ($\geq \$50$k). Restricting the index to the
$15$ markets whose YES outcome is unambiguously ``US recession occurs''
yields a striking exception: across the $18$-month window in which these
markets were actively traded, the recession-only pessimism index
co-moves strongly with ICS in first differences
($r_{\Delta} = -0.68,\ p = 0.003$;
OLS slope $-35.3$ (HAC s.e.\ $6.1$), $R^2 = 0.46$,
$p < 10^{-8}$). Granger-causality tests are null in both directions
and the lead--lag profile peaks at the contemporaneous lag,
so the relationship is informational co-movement rather than a forecast.
A placebo random walk shows a strongly significant levels correlation
($r = -0.56$, $p = 2\!\times\!10^{-4}$) but no relationship in differences,
underscoring that levels regressions on these series are unreliable. We
release the data-fetching and analysis pipelines for replication.\footnote{
Code and data sources: \url{https://github.com/<redacted>/papers/polymarket-vs-michigan-consumer-sentiment}.}
\end{abstract}

\section{Introduction}

Polymarket grew from $\sim$\$8M in monthly volume at the start of 2024
to over $\$2.5$ billion by November 2024 \citep{tsang2026anatomy}, an
expansion that has prompted both popular interest in the platform as a
``real-time sentiment thermometer'' and academic interest in whether
blockchain-based prediction markets offer comparable accuracy to traditional
forecasting tools \citep{cutting2025betting}. The University of Michigan
Index of Consumer Sentiment (ICS) is, in turn, one of the most-cited
survey-based measures of household economic perceptions and a documented
predictor of personal consumption growth \citep{carroll1994does,
ludvigson2004consumer, bram1998does}.

This paper asks a deliberately narrow question:
\begin{quote}
Does a liquidity-weighted index built from Polymarket's macro contracts
co-move with, or predict, the monthly Michigan ICS during 2022--2025?
\end{quote}

The framing is motivated by two distinct hypotheses one might hold. First,
a \emph{forecasting} hypothesis: because Polymarket prices update
continuously, they may aggregate news faster than monthly survey responses,
and so should lead the ICS. Second, a \emph{sentiment-mirror} hypothesis:
even if the markets carry no informational lead, an aggregate of
macro-pessimism contracts should at minimum track ICS contemporaneously,
because both should respond to the same macroeconomic news.

We find essentially no support for the aggregate version of either
hypothesis. A monthly index built from all $573$ macro-relevant markets
with non-trivial liquidity is uncorrelated with ICS at every horizon we
test. We do find a striking exception when we narrow attention to the
small set of markets explicitly resolving on a US recession declaration:
during the $18$ months in which these contracts traded actively
(2024-08 to 2026-01), the contemporaneous month-over-month
correlation between the recession-pessimism index and Michigan ICS is
$-0.68$, with an OLS $R^2$ of $0.46$. Even there, however,
Granger-causality tests offer no evidence that the market leads the
survey -- the relationship is a real-time mirror, not a forecast.

We frame our results not as a refutation of the value of prediction
markets but as a useful negative result that constrains the strong form
of the ``Polymarket as macro signal'' claim and motivates more targeted,
construct-validated index design.

\section{Related Work}

\paragraph{Prediction markets as forecasts.}
A long line of work documents that prediction market prices are
calibrated and informationally efficient under reasonable conditions
\citep{wolfers2004prediction, berg2008long, page2013do,
arrow2008promise, snowberg2013prediction}. Subject-by-subject
comparisons against experts and polls have generally favoured markets
when liquidity is non-trivial and outcomes are well-defined
\citep{berg2008long, atanasov2017distilling}. \cite{manski2006interpreting}
gives a precise treatment of when contract prices may or may not be
interpreted as the central tendency of beliefs.

\paragraph{Consumer sentiment as an economic indicator.}
The Michigan ICS and the comparable Conference Board index have been
shown to contain incremental predictive information for household
expenditure, beyond what is in observable macro data, although the
forecasting power is modest and concentrated at short horizons
\citep{carroll1994does, ludvigson2004consumer, bram1998does}. The
literature has therefore treated survey-based sentiment as a real but
noisy signal of household economic perceptions.

\paragraph{Bridging prediction markets and macroeconomic forecasting.}
\cite{snowberg2013prediction} survey the use of prediction-market prices
as macroeconomic forecasts, including the historic Iowa Electronic
Markets and the Hewlett-Packard internal markets. \cite{baker2016measuring}
construct a text-based Economic Policy Uncertainty (EPU) index that has
become a standard alternative signal to which new indices are often
compared. To our knowledge no published paper has constructed a
liquidity-weighted Polymarket index of macro pessimism and benchmarked it
against Michigan ICS at the monthly frequency, although working papers on
the 2024 election cycle \citep{tsang2026anatomy, cutting2025betting} have
begun to characterise Polymarket as a venue for macro and political
forecasting. \cite{servan2004does} provides early evidence that real-money
markets outperform play-money ones, and \cite{tetlock2007giving} is a
methodological precedent for relating high-frequency price-derived
sentiment indices to slower-moving surveys.

\section{Data}

We use three primary data sources, all free and public:

\paragraph{University of Michigan Index of Consumer Sentiment.}
We download the canonical \texttt{tbmics.csv} file directly from the
University of Michigan Survey Research Center
(\url{http://www.sca.isr.umich.edu/files/tbmics.csv}); this gives
monthly ICS values from November 1952 to April 2026 ($672$ rows). We use
the period overlapping our Polymarket coverage (Feb 2023 -- Apr 2026,
$39$ months).

\paragraph{Polymarket markets metadata.}
We page through the Polymarket Gamma REST API at
\url{https://gamma-api.polymarket.com/markets} using all eleven
macroeconomic-relevant tags discoverable via the \texttt{/tags} endpoint
(``recession'', ``Economy'', ``Macro Inflation'', ``Macro Indicators'',
``Fed Chair'', ``interest rates'', ``GDP'', ``economy1'',
``Macro Graph'', ``Macro Single'', ``initial jobless claims''),
supplemented by a slug-keyword filter over closed markets in four
overlapping date windows from 2022-01-01 to 2026-06-30. After
deduplication and a minimum-volume filter of \$5{,}000 we retain
$2{,}553$ candidate markets.

\paragraph{Polymarket daily price histories.}
For each candidate market we issue a request to the CLOB endpoint
\url{https://clob.polymarket.com/prices-history} with daily fidelity
(\texttt{interval=max}, \texttt{fidelity=1440}). We restrict to the
top $1{,}200$ markets by volume to keep the data pull tractable and
require at least ten daily price ticks per market. The pulled dataset
contains $1{,}107$ markets with usable history and $80{,}781$ daily
price ticks spanning 2023-02-03 to 2026-05-16.

\section{Methodology}

\subsection{Polarity classification}
\label{sec:polarity}
Each Polymarket contract has two outcomes labelled YES and NO with prices
$p_t \in [0,1]$ summing to $1$ at every instant up to market-maker
spread. We map each market to a polarity
$s_i \in \{-1, 0, +1\}$ based on its slug:
\begin{itemize}
\item $s_i = +1$ when the YES outcome corresponds to an unambiguously
  adverse macro event (e.g., \emph{will-the-next-recession-in-the-us-happen-by-q3-2022},
  \emph{us-recession-in-2025}, \emph{will-the-fed-increase-interest-rates-by-25-bps},
  \emph{will-jobless-claims-exceed-255k}, \emph{will-annual-inflation-increase-by-3pt0-or-more});
\item $s_i = -1$ when YES corresponds to a positive macro event (e.g.,
  \emph{fed-decreases-interest-rates}, \emph{will-us-gdp-growth-be-greater-than-2});
\item $s_i = 0$ (dropped) when YES has no clear macro interpretation.
\end{itemize}
The classification is implemented by a list of keyword patterns and a
small set of inflation-level regexes that handle ``or-more''/``or-less''
qualifiers (see \texttt{scripts/analysis.py:polarity}). It is admittedly
imperfect; we therefore report several subset results so that readers
can judge sensitivity to the labelling. Of $2{,}553$ candidate markets,
$573$ pass the polarity filter ($325$ with $s_i = +1$, $248$ with
$s_i = -1$) -- the remainder are either non-macro contracts that share
tags with macro ones (e.g., questions about specific Powell utterances)
or markets whose macro polarity is genuinely ambiguous.

\subsection{Pessimism index}
For each labelled market $i$ and day $t$, define a pessimism contribution
\[
\pi_{i,t} = \begin{cases} p_{i,t}, & s_i = +1, \\ 1 - p_{i,t}, & s_i = -1. \end{cases}
\]
The liquidity-weighted index on day $t$ is
\[
P_t = \frac{\sum_{i \in \mathcal{A}(t)} w_i\,\pi_{i,t}}{\sum_{i \in \mathcal{A}(t)} w_i},
\quad w_i = \min(\mathrm{vol}_i,\, q_{95}),
\]
where $\mathcal{A}(t)$ is the set of markets with a price tick on day
$t$ and $q_{95}$ is the $95$th percentile of cumulative volume across
the kept set; the cap prevents a single $\$10^8$-volume market from
dominating the index. We resample $P_t$ to monthly means
$P_{m}$ before merging with the monthly Michigan ICS.

\subsection{Statistical tests}
Because both series are likely non-stationary (we report ADF $p$-values
in Table~\ref{tab:headline}), our preferred specification is the
first-difference regression
\[
\Delta \mathrm{ICS}_m = \alpha + \beta\,\Delta P_m + \varepsilon_m,
\]
estimated by OLS with Newey--West heteroscedasticity- and
autocorrelation-consistent standard errors using $3$ lags
\citep{newey1987simple}. We also report Pearson and Spearman correlations
on both levels and differences, and Granger-causality tests on first
differences at lags $1,2,3$. A lead--lag cross-correlation scan
($\tau \in \{-3, \ldots, +3\}$ months) is reported to detect any
forecasting power masked by zero-lag specifications.

\paragraph{Placebo.} As a sanity check on the levels-regression results
we replace the Polymarket index by a Brownian random walk of the same
length and re-run the full battery. We pre-register the prediction
that a properly differenced design will null the placebo.

\section{Experiments}

\subsection{Sample composition}

Table~\ref{tab:composition} summarises the subsets used in the
analysis. The ``full'' index aggregates all $573$ labelled markets;
subsets restrict to particular semantic families. ``\#markets'' is the
count passing both the polarity filter and the \$5{,}000 minimum-volume
filter; ``high-vol share'' counts those with reported volume above
\$1\,M; ``\#months\,$\cap$\,ICS'' is the number of months where both
the subset's monthly mean and the Michigan ICS are defined.

\begin{table}[h]
\centering
\caption{Composition of the seven indices we evaluate.}
\label{tab:composition}
\begin{tabular}{lrrrl}
\toprule
Index & \#markets & \#months\,$\cap$\,ICS & high-vol share & period \\
\midrule
Full liquidity-weighted    & 573 & 39 & 59  & 2023-02--2026-05 \\
Full equal-weighted        & 573 & 39 & 59  & 2023-02--2026-05 \\
Recession only             & 15  & 18 & 1   & 2024-08--2026-01 \\
Fed only                   & 138 & 39 & 57  & 2023-02--2026-05 \\
Inflation only             & 268 & 22 & 1   & 2023-12--2026-04 \\
High volume only ($\geq \$50$k) & 293 & 39 & 59  & 2023-02--2026-05 \\
Placebo random walk        & --  & 39 & --  & 2023-02--2026-05 \\
\bottomrule
\end{tabular}
\end{table}

\subsection{Visual inspection}

Figure~\ref{fig:figure1}(a) overlays the daily liquidity-weighted
pessimism index against Michigan ICS. ICS falls from approximately
$72$ in early 2024 to under $52$ by April 2025 -- a $\sim 28\%$
collapse driven by the post-election tariff and tax-policy uncertainty
shock that is documented in contemporaneous press coverage and the
EPU index \citep{baker2016measuring}. The Polymarket pessimism index,
in contrast, oscillates between $0.3$ and $0.7$ over the same window
with no comparable trend; visual eyeball-correlation is weak at best.

\begin{figure}[t]
\centering
\includegraphics[width=\linewidth]{figure_1.png}
\caption{(a) Time series of monthly Michigan ICS (blue, left axis) and
daily liquidity-weighted Polymarket pessimism index (red, right axis).
(b) Scatter of monthly first differences using the full index;
the OLS slope is statistically indistinguishable from zero.
(c) The same plot restricted to the recession-only subset of $15$
markets: a sharp negative slope appears
($r_{\Delta} = -0.68$, $n = 17$).}
\label{fig:figure1}
\end{figure}

\section{Results}

\subsection{Main null result}

The full liquidity-weighted index is uncorrelated with Michigan ICS
in both levels and first differences (Table~\ref{tab:headline}, row
``Full''). The OLS first-difference coefficient is $-1.2$
($t \approx -0.14$); $R^2 = 0.0007$. We cannot reject the null of
zero relationship under any standard test, and the absence of signal
holds when we restrict to higher-volume markets only or use
equal-weighting. The equal-weighted variant is marginally
\emph{positively} associated with ICS in differences
($\beta = +13.2$, $p = 0.007$), the opposite of what a
``pessimism rises $\Rightarrow$ sentiment falls'' theory predicts;
we interpret this as a noise pattern rather than a real reversal,
since it is driven by a handful of high-leverage observations in
the political/general-economy markets where polarity assignment is
particularly contestable.

\begin{table}[t]
\centering
\caption{Headline regression and correlation results. All $p$-values
are HAC (Newey--West, $3$ lags) for the OLS row and standard for
correlations. Sample sizes $n$ are the number of months of overlap
with the Michigan ICS series.}
\label{tab:headline}
\small
\begin{tabular}{lrrrrrr}
\toprule
Index & $n$ & $r_{\Delta}$ & $p$ & $\beta_{\Delta}$ (HAC) & $p$ & $R^2_{\Delta}$ \\
\midrule
Full liquidity-weighted & 39 & $-0.027$ & $0.87$ & $-1.23\ (8.6)$ & $0.89$ & $0.001$ \\
Full equal-weighted     & 39 & $+0.318$ & $0.05$ & $+13.20\ (4.9)$ & $0.007$ & $0.101$ \\
Recession only          & 18 & $\mathbf{-0.675}$ & $\mathbf{0.003}$ & $\mathbf{-35.32\ (6.1)}$ & $\mathbf{<10^{-8}}$ & $\mathbf{0.456}$ \\
Fed only                & 39 & $-0.012$ & $0.94$ & $-0.52\ (8.3)$  & $0.95$ & $0.000$ \\
Inflation only          & 22 & $+0.385$ & $0.08$ & $+10.66\ (4.1)$ & $0.010$ & $0.148$ \\
High volume ($\geq\$50$k) & 39 & $-0.012$ & $0.94$ & $-0.51\ (8.3)$  & $0.95$ & $0.000$ \\
\midrule
Placebo random walk     & 39 & $+0.205$ & $0.22$ & $+6.63\ (8.2)$  & $0.42$ & $0.042$ \\
\bottomrule
\end{tabular}
\end{table}

ADF $p$-values for the full pessimism index are $0.37$ (levels) and
the ICS series gives $0.97$, so in neither case can we reject a unit
root in levels. This is why we treat the first-difference results as
the load-bearing tests and report the levels results only for
completeness.

\subsection{Recession subset is the exception}

The recession-only subset shows a striking pattern.
Monthly first differences in the recession-pessimism index and
Michigan ICS have correlation $-0.675$ ($p = 0.003$, $n = 17$ paired
months). The Newey--West OLS slope is $-35.3$ (standard error $6.1$,
HAC $p < 10^{-8}$) and $R^2 = 0.46$. Mechanically: a $10$-percentage-point
rise in the index-implied recession probability over the month
co-occurs with a $\sim 3.5$-point fall in Michigan ICS.
Figure~\ref{fig:figure1}(c) shows the scatter.

The fact that the result is driven by recession markets -- and only by
those -- has an intuitive reading: these contracts have the cleanest
\emph{construct validity} as a measure of household-relevant
macro pessimism, since they resolve on the exact event most consumer
sentiment surveys ask about. By contrast, Fed-rate contracts encode
trading hypotheses about the central bank's reaction function, and
inflation-level contracts price short-window CPI realisations
that may move independently of household perceptions.

\subsection{Lead--lag and Granger tests are null}

We test the stronger forecasting claim by (i) running a Granger-causality
F-test on first differences at lags $1,2,3$ for the recession-pessimism
$\to$ ICS direction, and (ii) computing the lead--lag cross-correlation
profile from $\tau = -3$ to $+3$ months
(Table~\ref{tab:leadlag}). Granger-causality tests reject neither the
``markets predict ICS'' nor the ``ICS predicts markets'' null at any
lag we examined. The lead--lag profile for the recession subset peaks
sharply at $\tau = 0$ ($r = -0.68$); the only other nominally
``significant'' lag at $5\%$ is $\tau = -1$ ($r = -0.47$, $p = 0.06$),
consistent with the small ICS-day-of-survey timing window. We
interpret the recession-subset finding as a real-time co-movement, not
a forecast.

\begin{table}[t]
\centering
\caption{Lead--lag cross-correlation for the recession-only subset.
Negative $\tau$ means ICS leads Polymarket by $|\tau|$ months;
positive $\tau$ means Polymarket leads ICS.}
\label{tab:leadlag}
\small
\begin{tabular}{lrrrrrrr}
\toprule
$\tau$ (months)  & $-3$ & $-2$ & $-1$ & $0$ & $+1$ & $+2$ & $+3$ \\
\midrule
$r_{\Delta}$ & $+0.30$ & $-0.21$ & $-0.47$ & $\mathbf{-0.68}$ & $-0.26$ & $+0.33$ & $+0.32$ \\
$p$          & $0.30$  & $0.46$  & $0.064$ & $\mathbf{0.003}$ & $0.34$  & $0.23$  & $0.27$  \\
$n$          & $14$    & $15$    & $16$    & $17$    & $16$    & $15$    & $14$    \\
\bottomrule
\end{tabular}
\end{table}

\subsection{Placebo test confirms differencing is necessary}

Substituting a Brownian random walk of length $39$ for the Polymarket
pessimism index produces a strongly significant \emph{levels}
correlation against the Michigan ICS ($r = -0.56$, $p = 2\!\times\!10^{-4}$,
$R^2 = 0.32$) -- a textbook spurious regression. The same placebo's
first-difference correlation is null ($r = +0.20$, $p = 0.22$),
matching the pre-registered prediction of the placebo paragraph above. We take this as direct empirical evidence that any
finding from levels regressions on this data should be treated with
suspicion, and we restrict our headline interpretation to the
first-difference results throughout.

\section{Discussion}

The full liquidity-weighted Polymarket pessimism index does not move
with the Michigan Index of Consumer Sentiment at the monthly frequency.
The result holds across a range of weighting and filtering choices, and
the only subset that shows a strong relationship is the narrow set of
contracts whose resolution criterion most closely matches what consumer
sentiment surveys ask about.

Two complementary explanations are consistent with our findings. First,
\emph{construct invalidity}: many of Polymarket's macro contracts are
designed for traders, not for measuring household sentiment. Markets on
``Will the Fed cut by 25 bps at the September meeting?'' encode a
finely-grained policy-path question whose price is dominated by
sophisticated traders' read of Fed Funds futures, not by the
consumer-relevant ``Is the economy headed for trouble?'' construct that
ICS captures. Aggregating these heterogeneous contracts dilutes the
signal we are after with orthogonal noise. Second, \emph{audience
mismatch}: Polymarket's user base is overwhelmingly composed of
politically engaged, finance-literate, crypto-native traders, whose
beliefs need not move with those of the broader US household population
the ICS samples.

The recession-subset result is consistent with both readings: a small
number of high-volume contracts that resolve directly on the
recession concept successfully reflect the same news events that move
household sentiment, but neither leads the other. From the perspective
of a researcher who wants a daily-frequency proxy for household
recession fears, the recession-only subset is a reasonable candidate;
from the perspective of an investor seeking a forecasting edge over the
ICS itself, our null Granger results give no encouragement.

\section{Limitations}

\paragraph{Short coverage window.}
The Polymarket macro-contract universe with non-trivial liquidity is
predominantly a 2024-onwards phenomenon. Our overlap with the ICS is
$39$ months for the full index and only $18$ months for the
recession subset, which is well within range where a $p < 0.05$ Granger
test could fail despite a real effect; we report no positive Granger
finding precisely because the data does not support one rather than
because we believe none exists.

\paragraph{Imperfect polarity labelling.}
Our slug-based polarity classifier labels only $573$ of $2{,}553$
candidate markets; the remainder are dropped. Mislabelling is possible
(particularly for inflation contracts at thresholds near the $\sim 3\%$
boundary we adopt); the inflation-subset result, which goes
``the wrong way'' for our hypothesis, may partly reflect this. We made
our classifier source-visible so that future replications can refine it.

\paragraph{Market-maker effects.}
Recent work \citep{tsang2026anatomy} shows that Polymarket reported
notional volume substantially overstates economically meaningful trade
volume once share creation/destruction is netted out. Our
liquidity weighting is based on the reported \texttt{volumeNum}
field and may therefore overweight markets whose true informational
volume is lower. We have not corrected for this; doing so requires
on-chain transaction-level data.

\paragraph{Single-country focus.}
The Michigan ICS is a US survey. We make no attempt to map non-US
Polymarket macro contracts (ECB, BoE, Argentina inflation, etc.) into
the index even though they are in the candidate set; we exclude them
when present.

\paragraph{Construct of ``macro pessimism''.}
We aggregate over heterogeneous events (recessions, Fed decisions,
inflation levels) with a $\pm 1$ polarity label. A more principled
construction would weight by an event's known elasticity to consumer
sentiment, which would require pre-existing structural estimates we did
not attempt to fit.

\section{Conclusion}

A liquidity-weighted Polymarket macro-pessimism index, naively
constructed from all macro-relevant resolved contracts between 2022 and
2025, does not track the Michigan Index of Consumer Sentiment in either
levels or first differences. A narrowly constructed sub-index using
only recession contracts \emph{does} co-move with ICS contemporaneously
in differences ($R^2 = 0.46$), but offers no out-of-sample forecasting
power over and above the survey itself. The result casts doubt on
informal claims that Polymarket is a useful aggregate macro-sentiment
signal at the monthly horizon while leaving open the possibility of
construct-validated, event-specific indices that could be informative.
All code, data sources, and intermediate artefacts are released in
the project repository so that future replications -- whether they
confirm or overturn these conclusions -- can extend the analysis
window and refine the polarity labelling.

\bibliographystyle{plainnat}
\bibliography{references}

\end{document}
